% =====================================================================
\section{本章小结}\label{sec:dy-summary}
% =====================================================================

本章的结论按严格性分级为：Born 截面、$\order{\as}$ 系数函数与 $\mu_F$ 无关性、单个割图中的 Glauber 钳夹、中间态取向决定 $k^-$ 极点、程函线间双胶子交换的 $\order{g^4}$ 幺正性在重整化微扰论中严格，单个割图给出指向未来的 Wilson 线与 DY、DIS 部分子分布的同一性由定义直接推出（\tagR）；旁观者 Glauber 相位的相消在阿贝尔程函模型中完整证明（\tagP）；区域、幺正性相消、同时形变、软因子为 $1$ 以及单举 DY、色单态末态、极化 DY 的共线因子化与 $1/Q^2$ 处的广义因子化给出全阶论证（\tagA），其精确形式是附录中以编号假设为前提的条件定理（\cref{tab:dy-correspondence}）；有色末态的因子化是猜想（\tagC），$m^4/s^2$ 上的广义因子化被破坏（\tagV）。

\begin{table}[htbp]
\centering
\caption{本章陈述与附录~O、附录~N 中精确形式的对应}
\label{tab:dy-correspondence}
\footnotesize
\setlength{\tabcolsep}{4pt}
\renewcommand{\arraystretch}{1.2}
\begin{tabularx}{\textwidth}{@{}>{\raggedright\arraybackslash}p{3.0cm}>{\raggedright\arraybackslash}X>{\raggedright\arraybackslash}p{5.2cm}@{}}
\toprule
本章陈述 & 精确形式 & 差别（假设、误差类型）\\
\midrule
\cref{lem:dy-fsi}（\tagA） & \cref{lem:appO-unitarity}（\tagR，$R_{\mathrm{cond}}$）；割位置恒等式\cref{lem:appO-cut-position} & 前提\cref{asm:appO-lcopt}显式化；在 E-S 型近似之下为精确恒等式\\
\cref{lem:dy-poles-orientation}（\tagR） & \cref{lem:appO-orientation}（\cref{sec:appO-orientation}） & 内容相同；附录给出复化后 $\Imag D$ 的下界及与\cref{lem:method-i0-compat}的对照\\
\cref{ex:dy-future-vs-past}（\tagR） & \cref{prop:appO-past-wilson}证明的第 (2) 步 & 单个割图给出未来指向的线；过去指向只在单侧改写之后\\
\cref{prop:dy-deformation}（\tagA） & \cref{prop:appO-one-sided}（E-W）、\cref{prop:appO-target}（精确几何）、\cref{prop:appO-eikonal}（E-P）、\cref{prop:appO-past-wilson}（\tagA） & 已证明的步骤与假设\cref{asm:appO-mixing}、\cref{asm:appO-deformation}、\cref{asm:appO-envelope}、\cref{asm:appO-rewritten-cs}分开\\
\cref{lem:dy-soft}（\tagA） & \cref{prop:appO-past-wilson}(b)、\cref{lem:method-soft-unitarity} & 退耦为 Wilson 线之后，软因子为 $1$ 是 $R_{\mathrm{def}}$ 的\\
\cref{prop:dy-abelian-model}（\tagP） & \cref{prop:appO-ss-cancel}（\tagR，$R_{\mathrm{cond}}$） & 推广到任意 $J_A$、$J_B$ 与交换族；只含非重合连接\\
\cref{prop:dy-eikonal-twogluon}（\tagR） & 证明见\cref{sec:appO-wilson}；参数化\cref{lem:appO-glauber-operator}；\cref{prop:appO-strict-window-shape} & 附录补充严格窗口对正规子形状的依赖\\
\cref{rem:dy-common-ordering}(iii) & \cref{prop:appO-twoexch-spectator}（\tagP，模型定理） & 复合旁观者、未割线、$\order{g^4}$；被割的直接线开放\\
\cref{thm:dy-factorization}（\tagA） & \cref{thm:appO-main}$\to$\cref{cor:appO-interface}$\to$\cref{prop:appN-dy}（均 \tagR，$R_{\mathrm{cond}}$） & 全部前提逐条列出；误差类型由 E-W 经接口化为 E-N\\
\bottomrule
\end{tabularx}
\end{table}

\begin{tcolorbox}[enhanced,breakable,title=本章涉及的常见误解,fonttitle=\bfseries,
  colback=gray!3!white,colframe=gray!55!black,boxrule=0.6pt,arc=2pt,
  left=5pt,right=5pt,top=3pt,bottom=3pt,before skip=6pt,after skip=6pt]
\small
\begin{itemize}[leftmargin=1.2em,itemsep=0pt,topsep=0pt,parsep=0pt]
\item “Glauber 胶子在单举截面中总是相消”：相消要求对旁观者的末态求和，并且测量在割线两侧作为恒等算符出现（\cref{lem:dy-fsi}的地位段落）；对 TMD 的 $pp\to h_1h_2X$ 与强子—强子硬衍射，它不成立（\cref{thm:break-tmd,prop:break-hh-diffraction}）。
\item “NNLO/N$^3$LO 与数据一致，说明 Glauber 相消已被检验”：部分子截面里没有旁观者（\cref{rem:dy-nnlo-not-glauber}）。
\item “把 $k^-$ 围道推到 $|k^-|\sim|\mathbf k_T|$ 就离开了 Glauber 区”：$|k^+k^-|$ 仍是 $\order{\lambda^3}$，须沿 $e_G$ 同时形变两个分量（\cref{prop:dy-deformation}）。
\item “非阿贝尔双 Glauber 交换留下非幺正残差”：那是严格截断的假象，精确的图给出 $\Real p_{\rm opp}=0$（\cref{prop:dy-eikonal-twogluon}，\cref{rem:dy-common-ordering}）。
\item “单个割图中的软 Wilson 线已经是指向过去的”：单个割图给出指向未来的线（\cref{ex:dy-future-vs-past}）。
\item “对 $q_T$ 微分的 DY，软因子也是 $1$”：只对共线单举截面成立；$q_T$ 谱的软因子进入 TMD 分布（\cref{lem:dy-soft}的地位段落）。
\item “把 $\mu_F$、$\mu_R$ 变化两倍得到的就是理论误差”：这是约定，因子化只保证残余依赖是高阶的（\cref{prop:dy-scale}之后的讨论）。
\item “余项是 $\Lambda^2/Q^2$，因为高扭度算符的维数高二”：这是 OPE 的陈述；对 DY，余项的幂次是猜想，重整化子分析只是证据（\cref{thm:dy-factorization}的地位段落，\cref{sec:dy-twist}）。
\item “所有强子—强子过程都像 DY 一样因子化”：色单态末态为 \tagA，有色末态只是 \tagC（\cref{thm:dy-colour-singlet}，\cref{conj:dy-coloured-final}）。
\end{itemize}
\end{tcolorbox}

后文的用途：阈值行为是重求和的出发点（\cref{prop:thr-oneloop}）；对 $q_T$ 微分的 DY，Glauber 分析给出指向过去的订书钉形 Wilson 线（\cref{thm:tmd-factorization,thm:tmd-sivers-sign}）；测量相容性判据与因子化的破坏见\cref{prop:break-measurement}，Glauber 有效理论见\cref{prop:front-syz}。
